% Appendix C, from exam/README.md: what the two written papers are for, what they cover, the
% conventions they assume, and how they are marked.

\chapter{The Self-Assessment Papers}
\label{exam:ch:about}

This appendix and the four after it hold two complete four-hour papers for \emph{UART: Hardware,
Driver \& Integration}, with worked solutions. Appendices~\ref{exa:ch:paper} and \ref{exb:ch:paper}
are the papers, questions only; Appendices~\ref{ansa:ch:answers} and \ref{ansb:ch:answers} are their
model answers, with marks.

\section*{What these are for}
The book has \textbf{no exam and nothing is marked}. What you get instead is the exercises after
every chapter, the provided self-checking testbench or host suite that grades each one, and the
questions in every chapter's review. Roughly half the code the course produces lives in those
exercises, and the testbenches in \file{hw/} and the suites in \file{fw/} are what say whether it
works.

These papers do not replace that. \textbf{They are here for you to test your own skills and
knowledge, nothing more.} They gate nothing, they are not a qualification, and no part of the book
requires them. Nothing in the companion repository depends on them: no module is built from them,
and neither \code{make build} nor \code{make format-check} knows they exist.

What they check is something the testbenches cannot. Every exercise is done at a keyboard, with the
chapter open beside you, GHDL ready to say no and \code{make test} ready to say no twice. These
papers ask what you can reconstruct with none of that in front of you --- which of the rules you
have absorbed, and which you have merely been looking up.

\textbf{Take one after the last chapter.} Every paper draws on all ten chapters and on both halves of
the system, so sitting one partway through examines material you have not read yet, and the result
says more about how far you have read than about what you have understood. The intended point is
after \chapref{10}, with the bench climbed and \code{app::EchoNode} echoing.

\section*{The two papers}
Both cover the whole book, Chapters~1 to~10, at the same weighting, in \textbf{eight questions: four
of VHDL, three of C++, and one with both halves on the bench at once.} Two of the eight span a pair
of chapters rather than one --- Question~3 covers the receive side across Chapters~3 and~4, and
Question~6 the driver and the suite that grades it across Chapters~7 and~8 --- which is why ten
chapters fit into eight questions without dropping anything. The papers share no question. Either
can be used alone; use both as a main sitting and a resit, or in alternate years.

\begin{booktable}{@{}p{5.2em}Lp{4.6em}p{2.6em}r@{}}
  \thead{Question} & \thead{Topic} & \thead{Chapters} & \thead{Side} & \thead{Marks} \\
  \midrule
  1 & The register map, the peripheral top, and binding & 1 & VHDL & 12 \\
  2 & Framing, the baud divider, and the transmitter & 2 & VHDL & 13 \\
  3 & The asynchronous line: synchronizer, oversampling, FIFO & 3, 4 & VHDL & 14 \\
  4 & The register bank: ports become registers & 5 & VHDL & 13 \\
  5 & The driver's contracts & 6 & C++ & 12 \\
  6 & The driver and the five-byte transaction & 7, 8 & C++ & 13 \\
  7 & The real transport and the freestanding target & 9 & C++ & 11 \\
  8 & Integration and bring-up & 10 & Both & 12 \\
  \midrule
    & & & & \textbf{100} \\
\end{booktable}

\textbf{Paper A leans towards reading and tracing.} Most of its questions hand you something --- a
\code{STATUS} word, a listing, a recorded byte log, a bench result --- and ask what it means: decode
a status read and say what the driver does next, trace a byte through the transmitter's frame
vector, work out how much baud mismatch the receiver's real sampling phase leaves, read five SPI
transactions off a stub's record and name the driver call that produced each.

Five of its parts do ask for code, so that ``leans towards'' does not mean ``never writes any'', and
they are chosen so that \textbf{no module is written on both papers}: \code{fifo}'s clocked process
and flags, the concurrent \code{STATUS} assembly, an AVR-portable header corrected, \file{env.cpp}'s
runtime symbols, and the \code{main()} that builds the stack on the target. Paper~B writes none of
those five.

\textbf{Paper B leans towards writing and reviewing.} Most of its questions ask you to produce a
module --- \code{baud_gen} in full, \code{sync} with its generic, the receiver's stop-bit branch, the
two \code{Interface} headers, \code{readReg} and \code{writeReg}, \code{AvrSpi},
\code{app::EchoNode} --- or to say what is wrong with one you are handed.

Both papers include at least one design that passes every check the course ships and is wrong
anyway, because that is the failure mode this subject actually has. Two more questions, one per
paper, are about the limits of the provided testbenches themselves, worked through a real gap that
the course's testbenches once had and no longer do.

\section*{Why the papers ask for both languages}
Because \textbf{the boundary is the subject}, and a boundary has two sides. A paper that examined
the VHDL and stopped at the register map would be examining digital design with VHDL; one that
started at \code{readReg} would be examining modern embedded C++. What this book adds is the thing
in between: one contract, written down once in Appendix~\ref{proto:ch}, implemented independently
on each side, and proven to agree only at the end.

So both papers cross the wire more than once. A question about the register bank's \code{STATUS}
bits ends up in the driver's \code{write()}; a question about byte order in \code{writeReg} ends up
in what \code{baud_gen} divides by. That is deliberate, and it is where the marks are.

What a written paper cannot do is run GHDL or \code{make test}, and it should not pretend to. The
rubric says outright that syntax is not what is being marked: a missing semicolon costs nothing, and
a process that describes a latch where the candidate meant a wire, or a \code{readReg} that
assembles its bytes backwards, costs that part in full.

\section*{Conventions the papers assume}
Both papers state these in their own rubric, so a candidate never has to have read this appendix.
\begin{itemize}
  \item \textbf{The protocol is the single source of truth.} Where an implementation and the
    protocol disagree, the implementation is wrong, and an answer that cites the protocol beats one
    that cites code.
  \item \textbf{VHDL-93}, with \code{ieee.std_logic_1164}, and \code{ieee.numeric_std} where a
    conversion needs it.
  \item \textbf{The VHDL conventions are the book's}: ports declared all inputs first, then all
    outputs; every \code{port map} and \code{generic map} positional; an active-low signal named
    \code{_n}; a two-flop synchronized signal named \code{_s2}; every submodule taking
    \code{reset_s2_n} while \code{uart_top} alone takes \code{reset_n}; SPI traffic between the two
    transport blocks prefixed \code{spi_}.
  \item \textbf{C++ is written in the book's AVR-portable subset} wherever it would run on the
    target: \code{<stdint.h>} and the bare types, no \code{std::}, nested \code{namespace} blocks
    rather than \code{namespace driver::uart}, no \code{[[nodiscard]]}, no \code{inline} variables.
    Host-only code, meaning the test suites, may use full modern C++, and the papers say which is
    which.
  \item \textbf{The C++ house style is expected}: \code{camelCase}, a leading capital on type names,
    \code{my} on private members and \code{our} on private statics, brace initialization,
    \code{noexcept} on driver methods, and copy and move deleted on anything holding a reference
    member or a unique hardware resource.
  \item \textbf{The numbers are fixed}: the DE0-CV runs at 50\,MHz (a 20\,ns period), the Nano at
    16\,MHz, \code{SCK} at f\textsubscript{osc}/16, which is 1\,MHz, the receiver oversamples 16
    times, and the default frame is 8N1.
  \item \textbf{Hardware answers, not code answers}, where a question asks what a design does. ``It
    assigns \code{x}'' is not an answer to ``what reaches the wire''.
\end{itemize}

Each paper supplies the clock rates, the \code{BAUD_DIV} formula and the SPI command-byte layout.
It does \textbf{not} supply the register map: Question~1 of Paper~A and Question~5 of Paper~B ask for part of it,
so it cannot be printed on the front page. Later questions consume it, which is what the follow-through
rule below is for.

\section*{Marking}
Every solution is written to be marked by somebody who has read the chapters and does not otherwise
write VHDL or freestanding C++ daily, so each carries the reasoning rather than the answer alone.
Marks are shown per part.

Four conventions are worth agreeing before a paper is marked:
\begin{itemize}
  \item \textbf{Method carries the marks.} A correct derivation with a slip in it is worth more than
    a correct answer with no working, and both papers are built so that later parts consume earlier
    ones. Follow an error through rather than penalising it twice: a driver that faithfully polls
    the wrong bit of a \code{STATUS} word the candidate decoded wrongly in Question~1 loses nothing
    in Question~6.
  \item \textbf{Mark the design, not the punctuation.} Where a VHDL listing is asked for, the marks
    are in the sensitivity list, the reset branch, the direction of the shift, whether the pulse is
    one clock wide, and whether every path assigns its output. Where C++ is asked for, they are in
    the byte order, the \code{const} and \code{noexcept} qualifiers that carry meaning, the ordering
    of the three register accesses, and whether the seam is respected.
  \item \textbf{The named traps are worth full marks on their own.} Several parts exist entirely to
    see whether a candidate avoids one specific mistake: a bit \emph{position} used as a mask, a
    \code{readReg} that assembles least significant byte first, an \code{RX_DATA} read that pops, a
    TX feeder that pops without watching \code{busy}, a \code{transfer()} with no \code{SPIF} poll, a
    receiver that leaves idle on a level rather than an edge, and a \code{run()} loop that blocks
    and can never be stopped. Where a solution flags one of these, a candidate who walks into it
    loses those marks and no others.
  \item \textbf{``State the consequence'' means the consequence.} Naming a defect earns half the
    marks; saying what it does to the running system --- which byte is lost, which poll never
    returns, what appears in the terminal --- earns the other half.
\end{itemize}

There is a fifth, particular to this subject. Several questions ask what a design does in
\textbf{simulation or on the host} and what it does \textbf{on the bench}, and the answer is that
they differ. Two questions go further and ask what a provided testbench \emph{cannot} catch. An
answer that treats a green testbench as proof has missed the point of the question, however correct
the rest of it is.
